\documentclass[12pt,a4paper]{article}
\usepackage[margin=2cm]{geometry}
\usepackage[spanish]{babel}
\usepackage{graphicx, multicol, latexsym, amsmath, amssymb}
\usepackage{array}
\usepackage{tabularx}
\usepackage{setspace}
\usepackage{titlesec}
\usepackage{subcaption}
\usepackage{float}
\usepackage{url}
\usepackage{hyperref}

% Ajuste de espacio entre secciones
\titlespacing*{\section}{0pt}{0.5em}{0.5em}

\renewcommand{\arraystretch}{1.3}

\begin{document}

% ---------------------------
% Encabezado con datos
% ---------------------------
\noindent
\begin{tabularx}{\textwidth}{|X|X|}
\hline
\textbf{Carné:} 2020045294 & \textbf{Nombre:} Justin Jaffeth Corea Masís \\
\hline
\textbf{Correo:} coreajustin288@estudiantec.cr & \textbf{Semana:} 6 del 17/09/2025 al 24/09/2025 \\
\hline
\end{tabularx}

\vspace{0.5cm}

% ---------------------------
% Tabla de actividades
% ---------------------------
\section*{Actividades realizadas}

% \begin{tabularx}{\textwidth}{|>{\raggedright\arraybackslash}p{8cm}|c|c|}
\begin{tabularx}{\textwidth}{|>{\raggedright\arraybackslash}p{12cm}|c|c|}
\hline
\textbf{Actividad} & \textbf{Fecha} & \textbf{Horas} \\
\hline
	Reunión con Paola para coordinar pruebas de campo & 19/09/2025 & 1 horas \\
\hline
	Probé diferentes aumentos para emular casos de placas cuyo reconocimiento falla & 19/09/2025 & 3 horas \\
\hline
	Mejoras y ampliación del marco teórico & 20/09/2025 & 8 horas \\
\hline
	Más pruebas de entrenamiento con aumentos más agresivos y variados & 22/09/2025 & 6 horas \\
\hline
	Más pruebas de entrenamiento con aumentos más agresivos y variados & 23/09/2025 & 6 horas \\
\hline
\end{tabularx}

\vspace{1cm}

% ---------------------------
% Firma
% ---------------------------
\noindent
\textbf{Firma del Profesor Asesor:}

\vspace{2cm}

\noindent\rule{8cm}{0.4pt}  
Nombre y firma

\newpage

\section{Resultados}

Alcance 98\% de precisión en validación, mejorando reconocimiento de placas que antes fallaban.
Algunos casos particulares siguen sin ser resueltos. 

\begin{figure}[H]
	\includegraphics[width=\textwidth]{/home/akai/Pictures/fail_real/exito_rate.png}
	\caption{Mejor resultado de entrenamiento}
\end{figure}

\begin{figure}[H]
	\begin{subfigure}[b]{0.3\textwidth}
		\includegraphics[width=\textwidth]{/home/akai/Pictures/fail_real/1.png}
	\end{subfigure}
	\begin{subfigure}[b]{0.3\textwidth}
		\includegraphics[width=\textwidth]{/home/akai/Pictures/fail_real/2.png}
	\end{subfigure}
	\begin{subfigure}[b]{0.3\textwidth}
		\includegraphics[width=\textwidth]{/home/akai/Pictures/fail_real/3.png}
	\end{subfigure}

	\begin{subfigure}[b]{0.3\textwidth}
		\includegraphics[width=\textwidth]{/home/akai/Pictures/fail_real/4.png}
	\end{subfigure}
	\begin{subfigure}[b]{0.3\textwidth}
		\includegraphics[width=\textwidth]{/home/akai/Pictures/fail_real/5.png}
	\end{subfigure}
	\begin{subfigure}[b]{0.3\textwidth}
		\includegraphics[width=\textwidth]{/home/akai/Pictures/fail_real/6.png}
	\end{subfigure}

	\begin{subfigure}[b]{0.3\textwidth}
		\includegraphics[width=\textwidth]{/home/akai/Pictures/fail_real/7.png}
	\end{subfigure}
	\begin{subfigure}[b]{0.3\textwidth}
		\includegraphics[width=\textwidth]{/home/akai/Pictures/fail_real/8.png}
	\end{subfigure}
	\begin{subfigure}[b]{0.3\textwidth}
		\includegraphics[width=\textwidth]{/home/akai/Pictures/fail_real/9.png}
	\end{subfigure}

	\begin{subfigure}[b]{0.45\textwidth}
		\includegraphics[width=\textwidth]{/home/akai/Pictures/fail_real/10.png}
	\end{subfigure}
	\begin{subfigure}[b]{0.45\textwidth}
		\includegraphics[width=\textwidth]{/home/akai/Pictures/fail_real/11.png}
	\end{subfigure}

	\caption{Casos de reconocimiento de placas que están fallando}
\end{figure}

\end{document}
